% !TeX root = ../初等数学笔记.tex
% 由主文档通过 \input 载入；请编译主文档。
\setcounter{chapter}{20} % 保留原书章号。
\chapter{复数及其几何意义}
\label{chap:em-complex}
方程 $x^2+1=0$ 在实数范围无解，促使我们扩展数系。复数扩展保留四则运算结构，并把乘法解释为平面上的缩放与旋转；它并不是给实数轴补上另一种大小顺序。

\section{复数、相等与四则运算}
\emph{复数}（complex number）写为 $z=x+iy$，其中 $x,y\in\mathbb R$，$i^2=-1$。也可严格地把它定义为实数对 $(x,y)$，规定加法为逐项相加、乘法为 $(x,y)(u,v)=(xu-yv,xv+yu)$，从而保证该表示和运算一致。实部为 $x$，虚部为实数 $y$，不是 $iy$。

复数相等当且仅当实部、虚部分别相等。展开并用 $i^2=-1$ 可得乘法公式；对 $u+iv\ne0$，
\[
 \frac{x+iy}{u+iv}=
 \frac{(xu+yv)+i(yu-xv)}{u^2+v^2}.
\]
分母为正，因此每个非零复数都有唯一倒数。不能给复数建立与实数有序域规则相容的顺序，否则平方都应非负，而 $i^2=-1<0$ 造成矛盾。

\section{共轭、模与复平面}
定义\emph{共轭}（complex conjugate）$\overline z=x-iy$ 和\emph{模}（modulus）$|z|=\sqrt{x^2+y^2}$。有
\[
 z\overline z=|z|^2,\qquad
 \overline{z+w}=\overline z+\overline w,\qquad
 \overline{zw}=\overline z\,\overline w.
\]
将乘积与其共轭相乘得 $|zw|^2=|z|^2|w|^2$，取非负平方根得到 $|zw|=|z||w|$。

在复平面中 $z$ 对应 $(x,y)$，加法是位移相加，共轭是关于实轴反射，$|z-w|$ 是两点距离。平面三角不等式给出 $|z+w|\le|z|+|w|$；含零值时取等，非零时取等要求二者同方向。$|z-a|=r$（$r>0$）表示以 $a$ 为圆心的圆，$|z-a|=|z-b|$（$a\ne b$）表示两点连线的垂直平分线。

\section{三角形式、辐角与乘除法}
若 $z\ne0$，记 $r=|z|>0$，则
\[
 z=r(\cos\theta+i\sin\theta).
\]
角 $\theta$ 称为\emph{辐角}（argument），所有可能值相差 $2k\pi$；需要单值时采用主辐角区间 $(-\pi,\pi]$。零复数没有确定方向，因此不定义辐角。

由和差角公式，两个非零复数相乘时模相乘、辐角相加；相除时模相除、辐角相减。乘以 $i$ 是逆时针旋转 $\pi/2$，乘以正数是等比例缩放，乘以负数再附加半周旋转。主辐角相加后可能落在规定区间外，须加减整周重新归入区间。

\section{复数乘方、开方与单位根}
对非零复数及整数 $n$，反复乘法和取倒数给出 de Moivre 公式
\[
 [r(\cos\theta+i\sin\theta)]^n
 =r^n(\cos n\theta+i\sin n\theta).
\]
设 $w=\rho(\cos\varphi+i\sin\varphi)\ne0$，正整数 $n$ 次方程 $z^n=w$ 的全部解为
\[
 z_k=\rho^{1/n}\left(\cos\frac{\varphi+2k\pi}{n}
       +i\sin\frac{\varphi+2k\pi}{n}\right),
 \qquad k=0,\ldots,n-1.
\]
必要性来自模的 $n$ 次方与辐角同余，充分性由代入验证；这 $n$ 个角模 $2\pi$ 互不相同，所以解也互异。$w=0$ 时唯一解为 $z=0$。

$z^n=1$ 的根称为 $n$ 次单位根，均匀分布在单位圆上。若 $\zeta\ne1$ 且 $\zeta^n=1$，则等比和给出 $1+\zeta+\cdots+\zeta^{n-1}=0$；这也反映了规则旋转配置的向量平衡。

\section{二次方程与几何应用}
实系数二次方程在 $\Delta<0$ 时，两根为
\[
 z=\frac{-b\pm i\sqrt{-\Delta}}{2a},\qquad a\ne0,
\]
互为共轭。更一般地，实系数多项式满足 $f(\overline z)=\overline{f(z)}$，故非实根成共轭对；这只说明已有根的对称性，不是一般多项式根存在性的证明。

围绕复点 $a$ 旋转 $\theta$ 的变换可写为
\[
 z'=a+(\cos\theta+i\sin\theta)(z-a).
\]
先减去中心、旋转位移、再加回中心，说明运算顺序的几何意义。复数坐标使旋转简洁，但长度关系仍需通过模表达，不能把复数本身当作有序长度。

\section{例题、习题与学习检查}
\begin{example}
求 $z^3=8$ 的全部复数解。
\end{example}
\begin{solve}
模为 $2$，辐角可取 $0,2\pi/3,4\pi/3$，故三根为 $2,-1+i\sqrt3,-1-i\sqrt3$。后三角表示中的整周重复不再产生新根；实数求根只保留了第一根。
\end{solve}
\begin{enumerate}
\item 化简 $(1+i)/(1-i)$。答：$i$。
\item 求 $z^2-2z+5=0$ 的根。答：$1\pm2i$。
\item 将 $z=2+i$ 绕 $a=1$ 逆时针旋转 $90^\circ$。答：$z'=1+i(1+i)=i$。
\end{enumerate}

